\chapter{Le funzioni}
\label{chap:funzioni}

Scomporre un problema in sotto-problemi più semplici è quasi sempre il modo
migliore per risolverlo (lo sviluppo top-down del
capitolo~\ref{chap:strutturata}), e per scrivere programmi grandi è vitale.
In C questo approccio si realizza con le \textbf{funzioni}. Nel
capitolo~\ref{chap:primi-programmi} ne abbiamo scritte alcune e abbiamo
incontrato prototipi e passaggio per valore; qui vediamo tutto in dettaglio:
come si definisce una funzione, come restituisce un risultato, che cosa
succede se il compilatore non la conosce, dove “vivono” le sue variabili.

\section{Programmi fatti di funzioni}

Un programma C è costituito da uno o più file (capitolo~\ref{chap:make}).
Ogni file contiene zero funzioni (come gli header) o più funzioni, e uno
solo contiene la funzione \code{main}, da cui comincia l'esecuzione;
\code{main} può chiamare altre funzioni, che a loro volta ne chiamano altre,
e tutte operano sulle variabili del programma.

Perché scrivere i programmi usando le funzioni?

\begin{itemize}
  \item è più facile scrivere \textbf{corrette} delle funzioni brevi;
  \item scrittura e ricerca degli errori (\textit{debugging}) sono più
        semplici, perché ogni funzione si prova da sola;
  \item le modifiche sono più semplici: si riscrivono solo le funzioni da
        cambiare;
  \item il codice è più comprensibile: funzioni brevi, con un nome che dice
        che cosa fanno, si leggono e si capiscono più facilmente.
\end{itemize}

\section{La definizione di una funzione}

\subsection{Intestazione e corpo}

Il codice che descrive che cosa fa una funzione è la sua
\term{definizione}. La forma generale è

\begin{lstlisting}[style=c]
type function_name(parameters)
{
    declarations
    statements
}
\end{lstlisting}

La parte prima della graffa aperta è l'\term{intestazione}
(\textit{header}); la parte fra le graffe è il \term{corpo} della funzione.
I \textbf{parametri} sono una lista di dichiarazioni di variabili separate da
virgole: sono detti anche \textbf{parametri formali}, perché ricevono i valori
(i \textit{parametri attuali}, o argomenti) solo al momento della chiamata.
La figura~\ref{fig:anatomia-funzione} mostra le parti nella funzione
fattoriale delle slide.

\begin{figure}[H]
  \centering
  \begin{tikzpicture}[font=\FiraCodeNL\small,
      pz/.style={anchor=base west, inner sep=0pt},
      lab/.style={font=\Lato\scriptsize, text=accent!70!black, align=center, inner sep=1pt}]
    \node[pz] (t) at (0,0) {int};
    \node[pz] (n) at ([xshift=5.4pt]t.base east) {factorial};
    \node[pz] (p) at (n.base east) {(int n)};
    \node[pz] (l2) at (0,-0.55) {\{};
    \node[pz] (l3) at (0,-1.1) {\ \ \ int i, product = 1;};
    \node[pz] (l4) at (0,-1.65) {\ \ \ for (i = 2; i <= n; i++)};
    \node[pz] (l5) at (0,-2.2) {\ \ \ \ \ \ product *= i;};
    \node[pz] (l6) at (0,-2.75) {\ \ \ return product;};
    \node[pz] (l7) at (0,-3.3) {\}};
    \draw[wtbrace] ([yshift=3pt]t.north east) -- ([yshift=3pt]t.north west);
    \draw[wtbrace] ([yshift=3pt]n.north east) -- ([yshift=3pt]n.north west);
    \draw[wtbrace] ([yshift=3pt]p.north east) -- ([yshift=3pt]p.north west);
    \node[lab, anchor=south] at ([yshift=0.32cm]t.north) {tipo del\\risultato};
    \node[lab, anchor=south] at ([yshift=0.32cm]n.north) {nome};
    \node[lab, anchor=south] at ([yshift=0.32cm]p.north) {parametro\\formale};
    \draw[wtbrace] (6.0,-0.12) -- (6.0,0.3);
    \node[lab, anchor=west] at (6.25,0.09) {intestazione (\textit{header})};
    \draw[wtbrace] (6.0,-3.4) -- (6.0,-0.4);
    \node[lab, anchor=west] at (6.25,-1.9) {corpo: dichiarazioni e istruzioni};
    \draw[wtarrowacc] (8.9,-2.75) -- ([xshift=4pt]l6.east);
    \node[lab, anchor=west] at (9.0,-2.75) {restituisce il risultato\\a chi ha chiamato};
  \end{tikzpicture}
  \caption{Le parti della definizione di una funzione.}
  \label{fig:anatomia-funzione}
\end{figure}

\code{factorial} riceve un \code{int} e restituisce un \code{int}: il
prodotto dei numeri da 2 a \code{n}, cioè $n!$ (per \code{n} uguale a 0 o 1
il ciclo non fa nessun giro, e il risultato è 1). Il valore restituito, se
necessario, viene convertito nel tipo della funzione, qui \code{int}; e
grazie alla lista dei parametri anche l'argomento, se necessario, viene
convertito in \code{int}.

\begin{affianco}[0.5][c]
\begin{lstlisting}[style=c, name=fattoriale.c]
#include <stdio.h>

int factorial(int n)              /* header */
{                                 /* body starts here */
   int i, product = 1;

   for (i = 2; i <= n; i++)
      product *= i;
   return product;
}

int main(void)
{
   int k;

   for (k = 0; k <= 5; k++)
      printf("%d! = %d\n", k, factorial(k));
   return 0;
}
\end{lstlisting}
\risultato
\begin{console}
$ ./fattoriale
0! = 1
1! = 1
2! = 2
3! = 6
4! = 24
5! = 120
\end{console}
\end{affianco}

Il corpo di una funzione è un blocco, e può contenere dichiarazioni di
variabili: \code{i} e \code{product} sono visibili \textbf{solo} dentro
\code{factorial}.

\subsection{\code{void}: niente risultato, niente parametri}

Una funzione può non restituire nulla, e può non avere parametri. In
entrambi i casi si usa la parola chiave \code{void} (“vuoto”):

\begin{affianco}[0.6][c]
\codicepiccolo
\begin{lstlisting}[style=c, name=santa.c]
#include <stdio.h>

void wrt_address(void)
{
  printf(" \n%s \n%s \n%s \n%s \n\n",
    "            ************************",
    "            **    SANTA CLAUS     **",
    "            **    NORTH POLE      **",
    "            **    EARTH           **");
}

int main(void)
{
  wrt_address();
  return 0;
}
\end{lstlisting}
\risultato
\begin{console}
$ ./santa

            ************************
            **    SANTA CLAUS     **
            **    NORTH POLE      **
            **    EARTH           **

\end{console}
\end{affianco}

Il primo \code{void} dice al compilatore che la funzione non restituisce
alcun valore; il secondo, fra le parentesi, che non ha parametri. La funzione
si chiama con l'istruzione \code{wrt\_address();}, con le parentesi vuote.
E se si scrivesse \code{a = wrt\_address();}? Non c'è nessun valore da
assegnare, e il compilatore si rifiuta:

\begin{console}
$ gcc -Wall -o santa2 santa2.c
santa2.c: In function (*@‘@*)main(*@’@*):
santa2.c:12:5: error: void value not ignored as it ought to be
   12 |   a = wrt_address();
      |     ^
\end{console}

\warning{Una svista nelle slide}{
  In una delle slide la funzione è scritta con \code{Printf}, con la P
  maiuscola. Il C distingue maiuscole e minuscole, e \code{Printf} non è
  \code{printf}: il compilatore lo tratta come una funzione sconosciuta
  (e suggerisce \textit{did you mean ‘printf’?}), poi il linker non la trova
  da nessuna parte e la compilazione fallisce con \textit{undefined reference
  to ‘Printf’}.}

\subsection{Altri esempi}

Ecco tre definizioni delle slide: una funzione che non fa nulla, una che
raddoppia un reale e una che somma tre interi.

\begin{affianco}[0.5][c]
\begin{lstlisting}[style=c, name=esempi.c]
#include <stdio.h>

void nothing(void)
{
}

double twice(double x)
{
   return (2.0 * x);
}

int all_add(int a, int b, int c)
{
   return (a + b + c);
}

int main(void)
{
   nothing();
   printf("%f\n", twice(3.5));
   printf("%d\n", all_add(1, 2, 3));
   return 0;
}
\end{lstlisting}
\risultato
\begin{console}
$ ./esempi
7.000000
6
\end{console}
\vspace{2pt}
\wtnota{\code{nothing} ha il corpo vuoto: chiamarla non fa nulla, ma è una
funzione lecita (può servire come segnaposto, in attesa di scriverla).}
\end{affianco}

Se il tipo della funzione non è indicato, il C tradizionale lo considera
\code{int}: \code{all\_add(int a, int b, int c) \{ \dots \}} equivale a
\code{int all\_add(\dots)}. Dal C99 questa regola non vale più, e
\code{gcc} protesta (\textit{return type defaults to ‘int’}); i compilatori
più recenti la considerano un errore. Specificare sempre il tipo della
funzione è comunque una buona pratica.

\subsection{Le definizioni nel C tradizionale}

Nel C tradizionale (quello precedente all'ANSI C) la definizione aveva una
sintassi diversa: fra le parentesi c'erano solo i nomi dei parametri, e le
loro dichiarazioni venivano dopo, prima della graffa aperta. L'ordine delle
dichiarazioni non contava, e senza parametri si mettevano solo le parentesi:

\begin{affianco}[0.5][c]
\begin{lstlisting}[style=c]
void f(a, b, c, x, y)
int     a, b, c;
double  x, y;
{
    ...
}
\end{lstlisting}
\risultato
\begin{lstlisting}[style=c]
void f()
{
    ...
}
\end{lstlisting}
\end{affianco}

Questa forma è obsoleta, e fra poco vedremo perché è pericolosa. Si
incontra ancora nei programmi vecchi.

\section{L'istruzione \code{return}}

\subsection{Terminare e restituire}

Con l'istruzione \code{return} la funzione termina e il controllo torna
all'ambiente da cui è stata chiamata. \code{return} non deve
necessariamente contenere un'espressione: \code{return;},
\code{return ++a;} e \code{return (a * b);} sono tutte istruzioni lecite (le
parentesi attorno all'espressione sono facoltative).

\begin{itemize}
  \item Se c'è un'espressione, il suo valore viene passato all'ambiente
        chiamante, convertito se necessario nel tipo della funzione.
  \item In una funzione ci possono essere zero o più \code{return}.
  \item Se non ce n'è nessuno, il controllo torna al chiamante quando si
        raggiunge la graffa di chiusura.
\end{itemize}

\begin{affianco}[0.5][c]
\begin{lstlisting}[style=c, name=assoluto.c]
#include <stdio.h>

double absolute_value(double x)
{
  if (x >= 0)
    return x;
  else
    return -x;
}

int main(void)
{
  printf("%f %f\n", absolute_value(-2.5),
                    absolute_value(4.0));
  return 0;
}
\end{lstlisting}
\risultato
\begin{console}
$ ./assoluto
2.500000 4.000000
\end{console}
\vspace{2pt}
\wtnota{Due \code{return}: a seconda del segno di \code{x}, la funzione
termina con l'uno o con l'altro.}
\end{affianco}

\subsection{Il valore si può ignorare}

Non sempre si usa il valore restituito. Nel frammento delle slide

\begin{lstlisting}[style=c]
while (...) {
    getchar();        /* il carattere letto viene scartato */
    c = getchar();
    ...
}
\end{lstlisting}

il primo \code{getchar()} serve solo a saltare un carattere. Lo facciamo
anche con \code{printf}, senza saperlo: anche \code{printf} restituisce un
valore, il numero di caratteri scritti, che di solito si ignora.

\begin{affianco}[0.5][c]
\begin{lstlisting}[style=c, name=ritorno.c]
#include <stdio.h>

int main(void)
{
   int n;

   n = printf("ciao\n");
   printf("caratteri scritti: %d\n", n);
   return 0;
}
\end{lstlisting}
\risultato
\begin{console}
$ ./ritorno
ciao
caratteri scritti: 5
\end{console}
\wtnota{I caratteri sono 5: \code{c}, \code{i}, \code{a}, \code{o} e l'a
capo.}
\end{affianco}

\section{I prototipi}

\subsection{Dichiarare prima di usare}

Le funzioni devono essere \textbf{dichiarate prima di essere usate}. L'ANSI C
fornisce per questo il \term{prototipo}, che informa il compilatore sul
numero e sul tipo dei parametri e sul tipo del valore restituito. Il
prototipo

\begin{lstlisting}[style=c]
double sqrt(double);
\end{lstlisting}

dice che \code{sqrt} è una funzione, che prende un solo argomento di tipo
\code{double} e che restituisce un valore di tipo \code{double}. La forma
generale è

\begin{lstlisting}[style=c]
type function_name(parameter type list);
\end{lstlisting}

Gli identificatori dei parametri sono facoltativi: \code{void f(char c, int
i);} e \code{void f(char, int);} sono equivalenti. Il compilatore non li
usa, ma servono da documentazione, soprattutto se i nomi sono significativi:
conviene metterli sempre.

Nel C tradizionale la dichiarazione non conteneva la lista dei parametri:
\code{double sqrt();}. I prototipi sono preferibili perché, conoscendo i
tipi dei parametri, il compilatore \textbf{converte} opportunamente i valori
passati alla funzione.

\subsection{Che cosa succede senza prototipo}

Consideriamo la chiamata \code{sqrt(a)}, con \code{a} intero e uguale a 4.
Con il prototipo il compilatore sa che \code{sqrt} vuole un \code{double}, e
converte 4 in 4.0; senza, non sa nulla, e il risultato non è corretto:

\begin{affianco}[0.5][t]
\begin{lstlisting}[style=c, name=prova.c]
#include <stdio.h>

double sqrt();

int main()
{
  int a = 4;

  printf("sqrt(4) = %f \n", sqrt(a));
}
\end{lstlisting}
\begin{console}
$ gcc -o prova prova.c -lm
$ ./prova
sqrt(4) = 0.000000
\end{console}
\risultato
\begin{lstlisting}[style=c, name=prova.c]
#include <stdio.h>

double sqrt(double);

int main()
{
  int a = 4;

  printf("sqrt(4) = %f \n", sqrt(a));
}
\end{lstlisting}
\begin{console}
$ gcc -o prova prova.c -lm
$ ./prova
sqrt(4) = 2.000000
\end{console}
\end{affianco}

Perché zero? Sui PC a 64 bit gli argomenti interi e quelli reali viaggiano in
due registri diversi del processore. Senza prototipo il valore 4 viene
consegnato come intero, nel registro degli interi; \code{sqrt} va a prendere
il suo \code{double} nell'altro registro, e ci trova un valore che non
c'entra (qui, per caso, 0). Con il prototipo il compilatore converte 4 in
4.0 e lo consegna dove \code{sqrt} lo cerca (figura~\ref{fig:registri}).

\begin{figure}[H]
  \centering
  \begin{tikzpicture}[font=\Lato\scriptsize,
      reg/.style={draw=primary!60!black, fill=primary!5!white, line width=0.7pt,
        minimum width=2.4cm, minimum height=0.55cm, font=\FiraCodeNL\scriptsize},
      vuoto/.style={reg, fill=white, draw=neutral!40!white},
      chi/.style={wtnode, font=\Lato\scriptsize, minimum width=1.9cm}]
    \node[font=\Lato\scriptsize\bfseries, anchor=west] at (-1.6,1.05) {senza prototipo};
    \node[chi] (m1) at (-0.6,0) {\code{main}};
    \node[reg] (ri1) at (3.6,0.35) {4};
    \node[reg] (rr1) at (3.6,-0.45) {0.0};
    \node[mnota, anchor=west] at ([xshift=3pt]ri1.east) {registro degli interi};
    \node[mnota, anchor=west] at ([xshift=3pt]rr1.east) {registro dei reali};
    \node[chi] (s1) at (9.4,0) {\code{sqrt}};
    \draw[wtarrow] (m1.east) -- node[above, mnota, pos=0.45] {consegna 4} (ri1.west);
    \draw[wtarrow, draw=merr] (rr1.south) |- node[pos=0.75, below, mnota, text=merr]
      {\code{sqrt} cerca qui il suo \code{double} e trova 0.0} (9.4,-1.0) -- (s1.south);
    \begin{scope}[yshift=-2.9cm]
      \node[font=\Lato\scriptsize\bfseries, anchor=west] at (-1.6,1.05) {con prototipo};
      \node[chi] (m2) at (-0.6,0) {\code{main}};
      \node[vuoto] (ri2) at (3.6,0.35) {};
      \node[reg] (rr2) at (3.6,-0.45) {4.0};
      \node[mnota, anchor=west] at ([xshift=3pt]ri2.east) {registro degli interi};
      \node[mnota, anchor=west] at ([xshift=3pt]rr2.east) {registro dei reali};
      \node[chi] (s2) at (9.4,0) {\code{sqrt}};
      \draw[wtarrow] (m2.east) -- node[below left, mnota, pos=0.55] {converte e consegna 4.0} (rr2.west);
      \draw[wtarrow, draw=mok] (rr2.south) |- node[pos=0.75, below, mnota, text=mok]
        {\code{sqrt} trova 4.0} (9.4,-1.0) -- (s2.south);
    \end{scope}
  \end{tikzpicture}
  \caption{Senza prototipo l'argomento viene consegnato nel posto sbagliato.}
  \label{fig:registri}
\end{figure}

\code{sqrt}, in realtà, è una funzione che \code{gcc} conosce bene: anche
senza \code{-Wall} avvisa che l'argomento è un \code{int} dove ci si
aspetta un \code{double}. Con \code{clang}, che la conosce altrettanto bene,
il primo programma stampa addirittura 2.000000. Con una funzione qualsiasi,
però, nessun compilatore può rimediare: senza prototipo il risultato non è
garantito.

\subsection{Come il compilatore scopre le funzioni}

Per il compilatore una funzione può essere dichiarata in tre modi: da una
\textbf{chiamata}, da una \textbf{definizione} o da un \textbf{prototipo}.

\begin{itemize}
  \item Se incontra una chiamata \code{f(x)} prima di qualunque altra
        dichiarazione, il compilatore assume la dichiarazione
        \code{int f();}: un risultato \code{int} e nessuna ipotesi sui
        parametri, che il programmatore deve passare del tipo giusto (dal
        C99 è un errore, che \code{gcc} segnala con un avviso, come abbiamo
        visto nel capitolo~\ref{chap:primi-programmi}).
  \item Una definizione in stile tradizionale dichiara solo il nome e il tipo
        del risultato, non i parametri.
  \item Una definizione in stile ANSI, come un prototipo, dichiara tutto.
\end{itemize}

Le slide lo mostrano con tre versioni della stessa funzione, che calcola la
radice quadrata del suo parametro, chiamata con un \code{int}, un
\code{double} e un \code{long double} che valgono tutti 4. Ecco la prima,
con la definizione tradizionale:

\begin{affianco}[0.5][c]
\begin{lstlisting}[style=c, name=kr1.c]
#include <stdio.h>
#include <math.h>

float f(x)
long double x;
{
  return sqrt(x);
}

int main()
{
  int a = 4;
  double b = 4;
  long double c = 4;

  printf("parametro intero: %f \n", f(a));
  printf("parametro double: %f \n", f(b));
  printf("parametro long double: %f \n", f(c));
}
\end{lstlisting}
\risultato
\begin{console}
$ ./kr1
parametro intero: 0.000000
parametro double: 0.000000
parametro long double: 2.000000
\end{console}
\end{affianco}

Il codice contiene la definizione di \code{f}, ma solo la riga
\code{float f(x)} vale come dichiarazione: non si assume nulla sui
parametri, che vanno passati del tipo giusto, perché il compilatore non può
fare nessuna conversione. Così solo la terza chiamata, che passa davvero un
\code{long double}, dà il risultato corretto.

\subsection{Un parametro senza tipo}

Nella seconda versione il tipo di \code{x} manca del tutto. Il compilatore
assume la dichiarazione \code{int x;} (e avvisa: \textit{type of ‘x’
defaults to ‘int’}); ora funziona solo la chiamata con l'intero:

\begin{affianco}[0.5][c]
\begin{lstlisting}[style=c, name=kr2.c]
#include <stdio.h>
#include <math.h>

float f(x)
{
  return sqrt(x);
}
\end{lstlisting}
\risultato
\begin{console}
$ ./kr2
parametro intero: 2.000000
parametro double: -nan
parametro long double: -nan
\end{console}
\wtnota{\code{main} è identico a quello di \code{kr1.c}.}
\end{affianco}

Nelle slide le due chiamate sbagliate stampano 0.000000; sul nostro
computer stampano \code{-nan} (\textit{not a number}: la radice di un
numero negativo). Sono valori presi a caso dalla memoria, e possono cambiare
da un computer all'altro.

\subsection{La definizione ANSI}

Nella terza versione l'intestazione è \code{float f(long double x)}: questa
volta la riga che fa da dichiarazione è un \textbf{prototipo}. Il
compilatore conosce la lista dei parametri, e può convertire l'\code{int} e
il \code{double} in \code{long double}:

\begin{affianco}[0.5][c]
\begin{lstlisting}[style=c, name=ansi.c]
#include <stdio.h>
#include <math.h>

float f(long double x)
{
  return sqrt(x);
}
\end{lstlisting}
\risultato
\begin{console}
$ ./ansi
parametro intero: 2.000000
parametro double: 2.000000
parametro long double: 2.000000
\end{console}
\end{affianco}

\begin{center}
\renewcommand{\arraystretch}{1.2}
\begin{tabular}{@{}l c c c@{}}
  \toprule
  \textbf{definizione di \code{f}} & \textbf{con un \code{int}} &
  \textbf{con un \code{double}} & \textbf{con un \code{long double}} \\
  \midrule
  \code{float f(x) long double x;} & \textcolor{merr}{0.000000} & \textcolor{merr}{0.000000} & \textcolor{mok}{2.000000} \\
  \code{float f(x)}                & \textcolor{mok}{2.000000}  & \textcolor{merr}{-nan}     & \textcolor{merr}{-nan} \\
  \code{float f(long double x)}    & \textcolor{mok}{2.000000}  & \textcolor{mok}{2.000000}  & \textcolor{mok}{2.000000} \\
  \bottomrule
\end{tabular}
\end{center}

In definitiva: la definizione della funzione in stile ANSI, oppure il suo
prototipo, deve comparire \textbf{prima} che la funzione sia usata. È anche
uno dei motivi per cui si includono gli header: contengono i prototipi delle
funzioni di libreria.

\section{Il passaggio per valore}

\subsection{Una copia per ogni argomento}

Gli argomenti delle funzioni sono passati \textbf{per valore}: ogni argomento
viene calcolato, e il suo valore viene usato localmente al posto del
parametro corrispondente. Se a una funzione si passa una variabile, il suo
valore nell'ambiente chiamante \textbf{non cambia}, anche se la funzione
modifica il valore che riceve (lo abbiamo visto nella
figura~\ref{fig:per-valore}). L'esempio delle slide:

\begin{affianco}[0.58][c]
\codicepiccolo
\begin{lstlisting}[style=c, name=compute_sum.c]
#include <stdio.h>

int compute_sum(int n)       /* sum the integers from 1 to n */
{
   int   sum = 0;

   for ( ; n > 0; --n)       /* stored value of n is changed */
      sum += n;
   return sum;
}

int main(void)
{
   int   n = 3, sum, compute_sum(int);

   printf("%d\n", n);
   sum = compute_sum(n);
   printf("%d\n", n);
   printf("%d\n", sum);
   return 0;
}
\end{lstlisting}
\risultato
\begin{console}
$ ./compute_sum
3
3
6
\end{console}
\wtnota{\code{compute\_sum(int)}, fra le dichiarazioni di \code{main}, è un
prototipo: lecito anche lì.}
\end{affianco}

\subsection{La pila delle chiamate}

Dove si trovano le variabili di una funzione? Ogni chiamata riceve una
propria zona di memoria (un \textit{record di attivazione}), che contiene i
parametri e le variabili locali; le zone si impilano una sull'altra, come
piatti, e ognuna viene tolta quando la sua funzione termina: è la
\term{pila delle chiamate} (\textit{call stack}). La
figura~\ref{fig:pila} la segue durante l'esecuzione di
\code{compute\_sum.c}.

\begin{figure}[H]
  \centering
  \begin{tikzpicture}[font=\Lato\scriptsize]
    \foreach \col/\testo in {0/{1. prima della chiamata}, 1/{2. inizio di \code{compute\_sum}},
                             2/{3. al \code{return}}, 3/{4. dopo il ritorno}}
      \node[font=\Lato\scriptsize\bfseries] at (\col*4.7+1.4,1.25) {\testo};
    % 1
    \variabile{a1}{(1.6,0.0)}{n}{3}
    \variabile{b1}{(1.6,-0.6)}{sum}{\mvuota}
    \scopo{2}{m1}{main}{(a1-nome) (a1) (b1-nome) (b1)}
    % 2
    \begin{scope}[xshift=4.7cm]
      \variabile{a2}{(1.6,0.0)}{n}{3}
      \variabile{b2}{(1.6,-0.6)}{sum}{\mvuota}
      \scopo{2}{m2}{main}{(a2-nome) (a2) (b2-nome) (b2)}
      \variabile[fill=accent!10!white]{c2}{(1.6,-1.95)}{n}{3}
      \variabile[fill=accent!10!white]{d2}{(1.6,-2.55)}{sum}{0}
      \scopo{3}{f2}{compute\_sum}{(c2-nome) (c2) (d2-nome) (d2)}
      \draw[mcerca] (a2.east) to[out=-10, in=10] node[wtlab, right, fill=none] {copia} (c2.east);
    \end{scope}
    % 3
    \begin{scope}[xshift=9.4cm]
      \variabile{a3}{(1.6,0.0)}{n}{3}
      \variabile{b3}{(1.6,-0.6)}{sum}{\mvuota}
      \scopo{2}{m3}{main}{(a3-nome) (a3) (b3-nome) (b3)}
      \variabile[fill=accent!10!white]{c3}{(1.6,-1.95)}{n}{0}
      \variabile[fill=accent!10!white]{d3}{(1.6,-2.55)}{sum}{6}
      \scopo{3}{f3}{compute\_sum}{(c3-nome) (c3) (d3-nome) (d3)}
      \draw[mcerca] (d3.east) to[out=10, in=-10] node[wtlab, right, fill=none] {6} (b3.east);
    \end{scope}
    % 4
    \begin{scope}[xshift=14.1cm]
      \variabile{a4}{(1.6,0.0)}{n}{3}
      \variabile{b4}{(1.6,-0.6)}{sum}{6}
      \scopo{2}{m4}{main}{(a4-nome) (a4) (b4-nome) (b4)}
      \node[mnota, align=center] at (1.4,-2.25) {la zona di \code{compute\_sum}\\è stata tolta dalla pila};
    \end{scope}
  \end{tikzpicture}
  \caption{La pila delle chiamate: \code{compute\_sum} ha la sua \code{n}
           (una copia) e la sua \code{sum}; modifica la propria \code{n}, e
           quando termina la sua zona scompare.}
  \label{fig:pila}
\end{figure}

\code{compute\_sum} porta la \textbf{sua} \code{n} da 3 a 0 mentre somma; la
\code{n} di \code{main} è un'altra variabile, e resta 3. Se si vuole che una
funzione modifichi una variabile del chiamante, ci sono due strade:
restituire il nuovo valore come risultato (e assegnarlo nel chiamante), oppure
passare l'\textbf{indirizzo} della variabile, cioè un puntatore
(capitolo~\ref{chap:puntatori}).

\section{Visibilità: variabili locali e globali}

\subsection{Locali e globali}

Le variabili dichiarate nel corpo di una funzione sono dette
\term{locali}: esistono e sono visibili solo lì. Le variabili dichiarate
fuori da ogni funzione sono dette \term{globali} (o esterne): sono visibili
da tutte le funzioni del file che vengono dopo la dichiarazione. Che cosa
stampa questo programma?

\begin{affianco}[0.5][c]
\begin{lstlisting}[style=c, name=globale.c]
#include <stdio.h>

int a = 33;

int all_add(int a, int b, int c)
{
  return (a+b+c);
}

int main(void)
{
  int b = 15, c = 34;
  printf("%d \n", all_add(a,b,c));
}
\end{lstlisting}
\risultato
\begin{console}
$ ./globale
82
\end{console}
\vspace{2pt}
\wtnota{In \code{main} non c'è una \code{a} locale: \code{a} è quella globale,
che vale 33, e $33 + 15 + 34 = 82$.}
\end{affianco}

Ora aggiungiamo una \code{a} locale in \code{main}, con valore 7:

\begin{affianco}[0.5][c]
\begin{lstlisting}[style=c, name=locale.c]
int main(void)
{
  int a=7, b = 15, c = 34;
  printf("%d \n", all_add(a,b,c));
}
\end{lstlisting}
\risultato
\begin{console}
$ ./locale
56
\end{console}
\vspace{2pt}
\wtnota{La \code{a} locale \textbf{nasconde} quella globale: $7 + 15 + 34 = 56$.}
\end{affianco}

Dentro \code{main}, il nome \code{a} indica la variabile locale, che
nasconde la globale; dentro \code{all\_add}, \code{a} è il parametro, che a
sua volta nasconde la globale (figura~\ref{fig:visibilita}). In ogni punto
del programma vale la dichiarazione \textbf{più vicina}.

\begin{figure}[H]
  \centering
  \begin{tikzpicture}[font=\Lato\scriptsize]
    % globale.c
    \node[font=\Lato\scriptsize\bfseries] at (2.4,2.4) {globale.c};
    \variabile{g1}{(1.0,1.55)}{a}{33}
    \variabile{mb1}{(1.0,0.2)}{b}{15}
    \variabile{mc1}{(1.0,-0.4)}{c}{34}
    \scopo{2}{mm1}{main}{(mb1-nome) (mb1) (mc1-nome) (mc1)}
    \variabile[fill=accent!10!white]{fa1}{(3.9,0.5)}{a}{33}
    \variabile[fill=accent!10!white]{fb1}{(3.9,-0.1)}{b}{15}
    \variabile[fill=accent!10!white]{fc1}{(3.9,-0.7)}{c}{34}
    \scopo{3}{ff1}{all\_add}{(fa1-nome) (fa1) (fb1-nome) (fb1) (fc1-nome) (fc1)}
    \scopo{1}{file1}{il file: variabili globali}{(g1-nome) (g1) (mm1) (ff1)}
    \draw[mcerca] (g1.east) to[out=0, in=90] (fa1.north);
    \node[mnota, text=mok] at (2.5,-1.75) {risultato: $33 + 15 + 34 = 82$};
    % locale.c
    \begin{scope}[xshift=7.6cm]
      \node[font=\Lato\scriptsize\bfseries] at (2.4,2.4) {locale.c};
      \variabile[draw=neutral!35!white, text=neutral!45!black]{g2}{(1.0,1.55)}{a}{33}
      \variabile{ma2}{(1.0,0.5)}{a}{7}
      \variabile{mb2}{(1.0,-0.1)}{b}{15}
      \variabile{mc2}{(1.0,-0.7)}{c}{34}
      \scopo{2}{mm2}{main}{(ma2-nome) (ma2) (mb2-nome) (mb2) (mc2-nome) (mc2)}
      \variabile[fill=accent!10!white]{fa2}{(3.9,0.5)}{a}{7}
      \variabile[fill=accent!10!white]{fb2}{(3.9,-0.1)}{b}{15}
      \variabile[fill=accent!10!white]{fc2}{(3.9,-0.7)}{c}{34}
      \scopo{3}{ff2}{all\_add}{(fa2-nome) (fa2) (fb2-nome) (fb2) (fc2-nome) (fc2)}
      \scopo{1}{file2}{il file: variabili globali}{(g2-nome) (g2) (mm2) (ff2)}
      \draw[mcerca] (ma2.east) -- (fa2-nome.west);
      \node[mnota, anchor=west, text=neutral!45!black] at ([xshift=4pt]g2.east) {nascosta in \code{main}};
      \node[mnota, text=mok] at (2.5,-1.75) {risultato: $7 + 15 + 34 = 56$};
    \end{scope}
  \end{tikzpicture}
  \caption{Da dove viene la \code{a} passata ad \code{all\_add}: a sinistra
           è la globale, a destra quella locale di \code{main}.}
  \label{fig:visibilita}
\end{figure}

\subsection{I blocchi}

Più in generale, un \textbf{blocco} è una parte di codice delimitata da
graffe, e può contenere dichiarazioni di variabili, visibili solo
all'interno del blocco. Il corpo di una funzione è un blocco, ma un blocco
può stare anche dentro un altro:

\begin{affianco}[0.5][c]
\begin{lstlisting}[style=c, name=blocchi.c]
#include <stdio.h>

int main(void)
{
   int a = 1;

   {
      int a = 2;              /* nasconde la a esterna */
      int b = 3;

      printf("dentro il blocco: a = %d, b = %d\n", a, b);
   }
   printf("fuori dal blocco: a = %d\n", a);
   return 0;
}
\end{lstlisting}
\risultato
\begin{console}
$ ./blocchi
dentro il blocco: a = 2, b = 3
fuori dal blocco: a = 1
\end{console}
\wtnota{Fuori dal blocco \code{b} non esiste più: usarla sarebbe un errore.
L'opzione \code{-Wshadow} di \code{gcc} segnala le variabili che ne
nascondono altre.}
\end{affianco}

\section{Le classi di memorizzazione}

Ogni variabile e ogni funzione del C ha due attributi: il \textbf{tipo} e la
\term{classe di memorizzazione} (\textit{storage class}), che stabilisce
\textbf{dove} e \textbf{per quanto tempo} esiste e \textbf{da dove} è
visibile. Le classi sono quattro, corrispondenti alle parole chiave
\code{auto}, \code{extern}, \code{register} e \code{static}.

\subsection{\code{auto}}

Le variabili dichiarate dentro un blocco sono \code{auto} per default.
Quando si entra nel blocco, il sistema alloca la memoria per le sue variabili
locali; quando se ne esce, la libera (è la pila della
figura~\ref{fig:pila}). Poiché è la classe più comune, la parola chiave
\code{auto} non si scrive quasi mai.

\subsection{\code{extern}}

Un modo semplice, ma sconsigliato, di passare informazioni fra blocchi e
funzioni sono le variabili esterne. Una variabile è esterna se è definita
fuori da ogni blocco, e riceve automaticamente la classe \code{extern}. La
parola chiave \code{extern}, invece, dice al compilatore di cercare la
variabile \textbf{da un'altra parte}. Un esempio su due file:

\begin{affianco}[0.5][t]
\begin{lstlisting}[style=c, name=file1.c]
#include <stdio.h>

int a = 1, b = 2, c = 3;        /* variabili esterne */
int f(void);

int main(void)
{
  printf("%3d\n", f());
  printf("%3d %3d %3d\n", a, b, c);
  return 0;
}
\end{lstlisting}
\risultato
\begin{lstlisting}[style=c, name=file2.c]
extern int a;              /* definita altrove */

int f(void)
{
  int b, c;                  /* variabili locali */

  a = b = c = 4;
  return (a + b + c);
}
\end{lstlisting}
\begin{console}
$ gcc -Wall -c file1.c
$ gcc -Wall -c file2.c
$ gcc -o esterne file1.o file2.o
$ ./esterne
 12
  4   2   3
\end{console}
\end{affianco}

\begin{itemize}
  \item I due file sono compilati separatamente
        (capitolo~\ref{chap:make}). In \code{file1.c}, \code{a}, \code{b} e
        \code{c} sono variabili esterne.
  \item Compilando \code{file2.c}, la riga \code{extern int a;} dice al
        compilatore che \code{a} è globale e definita in un altro file.
  \item Il linker mette insieme le informazioni e fa sì che \code{f} usi la
        cella di memoria giusta (figura~\ref{fig:extern}).
  \item \code{f} assegna 4 alla \code{a} globale e alle sue \code{b} e
        \code{c} locali, e restituisce 12; le \code{b} e \code{c} globali
        restano 2 e 3.
\end{itemize}

\begin{figure}[H]
  \centering
  \begin{tikzpicture}[font=\FiraCodeNL\scriptsize,
      sim/.style={anchor=west, inner sep=1.5pt},
      box/.style={draw=primary!65!black, fill=primary!4!white, line width=0.7pt,
        rounded corners=3pt}]
    \draw[box] (0,0) rectangle (3.2,-2.6);
    \node[font=\Lato\scriptsize\bfseries, anchor=south west] at (0,0.02) {file1.o};
    \node[sim, text=mok] (da) at (0.1,-0.35) {D a};
    \node[sim, text=mok] at (0.1,-0.75) {D b};
    \node[sim, text=mok] at (0.1,-1.15) {D c};
    \node[sim] at (0.1,-1.55) {T main};
    \node[sim, text=merr] (uf) at (0.1,-1.95) {U f};
    \node[sim, text=neutral!45!black] at (0.1,-2.35) {U printf};
    \draw[box] (5.6,-0.6) rectangle (8.4,-1.8);
    \node[font=\Lato\scriptsize\bfseries, anchor=south west] at (5.6,-0.58) {file2.o};
    \node[sim, text=merr] (ua) at (5.7,-0.95) {U a};
    \node[sim, text=mok] (tf) at (5.7,-1.45) {T f};
    \draw[wtarrowacc] (ua.west) -- (da.east);
    \draw[wtarrowacc] (uf.east) -- (tf.west);
    \node[mnota, anchor=north, align=center] at (4.2,-2.75)
      {\code{nm}: \code{D} è una variabile definita nel file, \code{T} una funzione definita,\\
       \code{U} un nome usato ma definito altrove: il linker collega ogni \code{U} alla sua definizione};
  \end{tikzpicture}
  \caption{Che cosa vede il linker: \code{file2.o} usa la \code{a} definita in
           \code{file1.o}, e \code{file1.o} usa la \code{f} definita in
           \code{file2.o}.}
  \label{fig:extern}
\end{figure}

Le informazioni si possono passare anche come parametri delle funzioni, ed
è la strada da preferire: bisogna cercare di evitare le variabili globali.
Con i parametri il codice è più chiaro e si capisce meglio dove i valori
vengono modificati; una variabile globale, invece, può essere cambiata da
qualunque funzione di qualunque file.

Anche le funzioni hanno per default la classe \code{extern}: sono visibili
ovunque nel programma, anche se è scritto su più file. Si può aggiungere
\code{extern} alla definizione o al prototipo
(\code{extern double sin(double);}), ma in generale non serve.

\subsection{\code{register}}

La classe \code{register} chiede al compilatore di tenere la variabile, se
possibile, in un \textbf{registro} del processore, la memoria più veloce
(capitolo~\ref{chap:calcolatore}). I registri sono pochi e servono anche al
sistema, quindi la richiesta non sempre è soddisfatta: in quel caso la
variabile diventa \code{auto}. Le variabili globali e le funzioni non
possono essere \code{register}.

\code{register} è un tentativo di velocizzare il programma, e le variabili
vanno scelte con cura: poche, usate molto spesso. I candidati tipici sono
le variabili di controllo dei cicli. Dichiarata all'inizio di \code{main},
però, la variabile occupa il registro per tutta l'esecuzione della
funzione; è più efficiente racchiudere dichiarazione e ciclo in un blocco:

\begin{affianco}[0.5][c]
\begin{lstlisting}[style=c]
int main(void)
{
    register int i;
    ...
    for (i = 0; i < LIMIT; i++) {
        ...
    }
}
\end{lstlisting}
\risultato
\begin{lstlisting}[style=c]
int main(void)
{
    ...
    {
        register int i;
        for (i = 0; i < LIMIT; i++) {
            ...
        }
    }
}
\end{lstlisting}
\end{affianco}

Una variabile \code{register} non ha un indirizzo in memoria, quindi non le
si può applicare \code{\&}: \textit{address of register variable ‘i’
requested}, risponde \code{gcc}. Oggi i compilatori scelgono da soli, e
molto bene, quali variabili tenere nei registri: la parola chiave
\code{register} è di fatto superata.

\subsection{\code{static} per le variabili locali}

Le dichiarazioni \code{static} hanno due usi distinti. Il primo è
\textbf{conservare il valore} di una variabile locale anche dopo l'uscita
dal blocco: una variabile locale \code{static} viene inizializzata una sola
volta, e non sta nella pila ma in una zona di memoria che esiste per tutta
l'esecuzione del programma.

\begin{affianco}[0.5][c]
\begin{lstlisting}[style=c, name=statica.c]
#include <stdio.h>

void f(void)
{
   static int cnt = 0;
   int volte = 0;

   ++cnt;
   ++volte;
   printf("cnt = %d   volte = %d\n", cnt, volte);
}

int main(void)
{
   f();
   f();
   f();
   return 0;
}
\end{lstlisting}
\risultato
\begin{console}
$ ./statica
cnt = 1   volte = 1
cnt = 2   volte = 1
cnt = 3   volte = 1
\end{console}
\end{affianco}

La prima volta \code{cnt} viene inizializzata a 0 e incrementata a 1. La
seconda volta che \code{f} viene chiamata, \code{cnt} \textbf{non} viene
inizializzata di nuovo: conserva il valore 1 e lo porta a 2. \code{volte},
che è \code{auto}, nasce e muore a ogni chiamata, e vale sempre 1
(figura~\ref{fig:statica}). Le slide usano questa idea per fare una cosa sì
e una no: \code{if (cnt \% 2 == 0)} distingue le chiamate pari da quelle
dispari.

\begin{figure}[H]
  \centering
  \begin{tikzpicture}[font=\Lato\scriptsize]
    \foreach \k in {1,2,3} {
      \node[font=\Lato\scriptsize\bfseries] at (\k*3.6-1.0,1.3) {chiamata \k};
      \variabile[fill=accent!10!white]{v\k}{(\k*3.6-0.8,-1.3)}{volte}{1}
      \scopo{3}{fr\k}{f}{(v\k-nome) (v\k)}
    }
    \variabile{c1}{(2.8,0.35)}{cnt}{1}
    \variabile{c2}{(6.4,0.35)}{cnt}{2}
    \variabile{c3}{(10.0,0.35)}{cnt}{3}
    \draw[wtarrowacc] (c1.east) -- node[above, mnota, text=accent!70!black] {$+1$} (c2-nome.west);
    \draw[wtarrowacc] (c2.east) -- node[above, mnota, text=accent!70!black] {$+1$} (c3-nome.west);
    \node[mnota, anchor=west, align=left] at (11.2,0.35) {\code{static}: la stessa variabile,\\fuori dalla pila, conserva il valore};
    \node[mnota, anchor=west, align=left] at (11.2,-1.3) {\code{auto}: una variabile nuova\\a ogni chiamata, che riparte da 0};
  \end{tikzpicture}
  \caption{\code{cnt} attraversa le chiamate; \code{volte} viene creata e
           distrutta ogni volta.}
  \label{fig:statica}
\end{figure}

\subsection{\code{static} per le variabili globali}

Il secondo uso riguarda le variabili esterne. Una variabile globale
\code{static} ha una visibilità più ristretta: è visibile solo dalle
funzioni del suo file definite dopo di lei, e non da quelle di altri file.
Tentare di raggiungerla con \code{extern} non serve. L'esempio delle slide è
una famiglia di generatori di numeri pseudo-casuali:

\begin{lstlisting}[style=c, name=random.c]
/* A family of pseudo random number generators. */

#define    INITIAL_SEED         17
#define    MULTIPLIER           25173
#define    INCREMENT            13849
#define    MODULUS              65536
#define    FLOATING_MODULUS     65536.0

static unsigned seed = INITIAL_SEED;   /* esterna, ma privata di questo file */

unsigned random(void)
{
   seed = (MULTIPLIER * seed + INCREMENT) % MODULUS;
   return seed;
}

double probability(void)
{
   seed = (MULTIPLIER * seed + INCREMENT) % MODULUS;
   return (seed / FLOATING_MODULUS);
}
\end{lstlisting}

\begin{itemize}
  \item \code{random()} produce un intero fra 0 e \code{MODULUS} (escluso),
        \code{probability()} un reale fra 0 e 1.
  \item Entrambe usano il vecchio valore di \code{seed} (il \textit{seme})
        per calcolare il nuovo, e lo aggiornano: ogni numero dipende dal
        precedente.
  \item Se \code{seed} venisse modificata per errore da un'altra parte del
        programma, la sequenza cambierebbe. Siccome è \code{static}, è
        privata di questo file e nessun altro modulo può toccarla.
\end{itemize}

Lo verifichiamo con un programma di prova, e poi con uno che prova a
“rubare” il seme:

\begin{affianco}[0.5][c]
\begin{lstlisting}[style=c, name=prova.c]
#include <stdio.h>

unsigned random(void);
double probability(void);

int main(void)
{
   int i;

   for (i = 0; i < 5; i++)
      printf("%u ", random());
   printf("\n%f\n", probability());
   return 0;
}
\end{lstlisting}
\risultato
\begin{console}
$ gcc -c random.c
$ gcc -o prova prova.c random.o
$ ./prova
48574 61999 40372 31453 39802
0.537521
$ nm random.o
000000000000002a T probability
0000000000000000 T random
0000000000000000 d seed
\end{console}
\end{affianco}

Il primo numero è proprio $(25173 \cdot 17 + 13849) \bmod 65536 = 48574$.
In \code{nm} il seme compare con la \code{d} \textbf{minuscola}: la
minuscola indica un simbolo locale al file, che il linker non mostra agli
altri. Infatti un programma che dichiara \code{extern unsigned seed;} e
prova ad assegnargli 0 non viene collegato:

\begin{affianco}[0.42][c]
\begin{lstlisting}[style=c, name=ruba.c]
#include <stdio.h>

extern unsigned seed;
unsigned random(void);

int main(void)
{
   seed = 0;
   printf("%u\n", random());
   return 0;
}
\end{lstlisting}
\risultato
Il compilatore accetta \code{ruba.c}: \code{extern} promette che \code{seed}
è definita da qualche parte. È il linker a non trovarla, perché in
\code{random.o} il seme non è visibile da fuori (le righe con
\textit{warning} sono avvisi secondari, conseguenza dello stesso problema).
\end{affianco}

\begin{console}
$ gcc -o ruba ruba.c random.o
/usr/bin/ld: /tmp/ccEe4NI7.o: warning: relocation against `seed' in read-only section `.text'
/usr/bin/ld: /tmp/ccEe4NI7.o: in function `main':
ruba.c:(.text+0xa): undefined reference to `seed'
/usr/bin/ld: warning: creating DT_TEXTREL in a PIE
collect2: error: ld returned 1 exit status
\end{console}

\subsection{\code{static} per le funzioni}

Anche una funzione può essere \code{static}: allora è visibile solo nel file
in cui è definita, e solo le funzioni dello stesso file possono chiamarla.
Serve a definire moduli con funzioni \textbf{private}. Nell'esempio delle
slide due file definiscono ciascuno la propria \code{sf}, \code{static}, e
non c'è alcun conflitto:

\begin{affianco}[0.5][t]
\begin{lstlisting}[style=c, name=main.c]
#include <stdio.h>

void a(void);
void f() { puts("main f"); }
static void sf() { puts("main sf"); }

void m() {
    f();
    sf();
}

int main() {
    m();
    a();
    return 0;
}
\end{lstlisting}
\risultato
\begin{lstlisting}[style=c, name=file.c]
#include <stdio.h>

void f(void);
static void sf() { puts("a sf"); }

void a() {
    f();
    sf();
}
\end{lstlisting}
\begin{console}
$ gcc -c file.c -o file.o
$ gcc -c main.c -o main.o
$ gcc -o main main.o file.o
$ ./main
main f
main sf
main f
a sf
\end{console}
\end{affianco}

\code{puts} scrive una stringa e va a capo. \code{main} chiama \code{m}, che
chiama \code{f} e la \code{sf} del suo file; poi chiama \code{a}, che sta in
\code{file.c}: \code{a} chiama \code{f}, che esiste solo in \code{main.c} ed
è visibile ovunque, e la \code{sf} del \textbf{proprio} file
(figura~\ref{fig:static-funzioni}).

\begin{figure}[H]
  \centering
  \begin{tikzpicture}[font=\FiraCodeNL\scriptsize,
      fn/.style={draw=primary!60!black, fill=white, rounded corners=2pt, line width=0.6pt,
        minimum width=1.5cm, minimum height=0.5cm},
      fs/.style={fn, draw=accent!65!black, fill=accent!10!white}]
    \node[fn] (main) at (0,0) {main};
    \node[fn] (m) at (0,-1.1) {m};
    \node[fn] (f) at (-1.3,-2.3) {f};
    \node[fs] (sf1) at (1.3,-2.3) {sf};
    \scopo{1}{fm}{main.c}{(main) (m) (f) (sf1)}
    \node[fn] (a) at (5.4,-1.1) {a};
    \node[fs] (sf2) at (5.4,-2.3) {sf};
    \scopo{1}{ff}{file.c}{(a) (sf2)}
    \draw[wtarrow] (main) -- (m);
    \draw[wtarrow] (m) -- (f);
    \draw[wtarrow] (m) -- (sf1);
    \draw[wtarrow] (main.east) -| (a.north);
    \draw[wtarrow] (a) -- (sf2);
    \draw[wtarrow, draw=accent!70!black, rounded corners=3pt] (a.west) -- (3.4,-1.1) |- (0.2,-3.15)
      node[pos=0.75, below, mnota, text=accent!70!black] {\code{a} chiama la \code{f} di \code{main.c}} -| (f.south);
    \node[mnota, anchor=west, align=left] at (7.0,-2.3) {\code{static}: visibile\\solo nel suo file};
  \end{tikzpicture}
  \caption{Chi chiama chi: \code{f} è visibile da entrambi i file, ciascuna
           \code{sf} solo dal proprio.}
  \label{fig:static-funzioni}
\end{figure}

\subsection{Riassunto}

\begin{center}
\renewcommand{\arraystretch}{1.2}
\begin{tabular}{@{}l>{\raggedright\arraybackslash}p{4.4cm}>{\raggedright\arraybackslash}p{5.2cm}>{\raggedright\arraybackslash}p{4.6cm}@{}}
  \toprule
  \textbf{classe} & \textbf{dove si dichiara} & \textbf{da dove è visibile} & \textbf{quanto vive} \\
  \midrule
  \code{auto}                 & in un blocco (è il default)       & nel blocco                          & dall'ingresso all'uscita dal blocco \\
  \code{register}             & in un blocco                      & nel blocco                          & come \code{auto} \\
  \code{static}               & in un blocco                      & nel blocco                          & per tutto il programma \\
  esterna (\code{extern})     & fuori dalle funzioni              & nel file; negli altri file con \code{extern} & per tutto il programma \\
  \code{static}               & fuori dalle funzioni              & solo nel suo file                   & per tutto il programma \\
  \bottomrule
\end{tabular}
\end{center}

\insintesi{
  Una funzione ha un'\textbf{intestazione} (tipo del risultato, nome,
  parametri formali) e un \textbf{corpo}; \code{void} indica che non
  restituisce nulla o che non ha parametri. \code{return} termina la funzione
  e restituisce un valore, convertito nel tipo della funzione; il valore si
  può ignorare. Una funzione va dichiarata \textbf{prima} dell'uso, con la
  definizione in stile ANSI o con un \textbf{prototipo}: senza, il
  compilatore non converte gli argomenti e i risultati sono sbagliati. Gli
  argomenti sono passati \textbf{per valore}: ogni chiamata ha la sua zona
  nella pila, con copie dei valori. Le variabili locali sono visibili solo
  nel loro blocco e nascondono quelle esterne con lo stesso nome. Le classi
  di memorizzazione stabiliscono durata e visibilità: \code{auto} (il
  default), \code{extern} (variabili definite altrove), \code{register}
  (superata), \code{static} (variabili locali che conservano il valore,
  variabili e funzioni private di un file).}

\section{Riferimenti}

\begin{itemize}
  \item \textbf{A. Kelley, I. Pohl}, \textit{A Book on C}: capitolo~5
        (funzioni e classi di memorizzazione).
  \item \textbf{Function declarations} e \textbf{Storage-class specifiers}
        --- cppreference.com. \\
        \urlx{https://en.cppreference.com/w/c/language/function_declaration} \\
        \urlx{https://en.cppreference.com/w/c/language/storage_duration}
\end{itemize}
